\chapter{Topological dynamics}





\section{A}

We will survey here some basic
concepts and results of topological dynamics, which we will need in this article(monograph). Recall that an action $(g,x)\in G\times X\mapsto g\cdot x\in X$ of a topological group $G$ on a topological space $X$ is {\it continuous} if it is continuous as a map from $G\times X$ into $X$.
We will consider continuous actions of (Hausdorff) topological groups $G$  on (non-$\emptyset$) compact, Hausdorff spaces $X$.
Actually we are primarily interested in metrizable topological groups $G$ and in fact only separable metrizable ones.
So although we will state in this survey several standard results for general topological groups, we will often give a sketch of an argument for the metrizable case only.
It should be kept in mind that if $G$ is metrizable (equivalently is Hausdorff and has a countable nbhd basis at the identity), then $G$ admits a right-invariant compatible metric $d_r$, which of course can always be taken to be bounded by 1, by replacing it, if necessary, by $\tfrac{d_r}{1+d_r}$.  See, \ref{thm:birkhoff-kakutani} for a detailed treatment.







Let $G$ be a topological group and $X$ a compact, Hausdorff space. We equip $H(X)$, the group of homeomorphisms
of $X$, with the compact-open topology, i.e., the topology with subbasis
$\{f\in H(X):f(K)\subseteq V\}$, with $K\subseteq X$ compact, $V\subseteq
    X$ open.


\begin{theorem}
    $H(X)$ is a topological group, and a continuous action of $G$
    on $X$ is simply a continuous homomorphism of $G$ into $H(X)$.
\end{theorem}

\begin{proof}
    Let $S(K, V) = \{f \in H(X) \mid f(K) \subseteq V\}$ denote the subbasic open sets of the compact-open topology, where $K \subseteq X$ is compact and $V \subseteq X$ is open. We must show that inversion and composition are continuous.

    Put $\iota(f) = f^{-1}$,
    let $S(K, V)$ be a subbasic open set. Its preimage under $\iota$ is:
    $$ \iota^{-1}(S(K, V)) = \{f \in H(X) \mid f^{-1}(K) \subseteq V\} $$
    Since $f$ is a bijection, the condition $f^{-1}(K) \subseteq V$ is equivalent to $f(X \setminus V) \subseteq X \setminus K$.

    Let $C = X \setminus V$ and $U = X \setminus K$. Because $X$ is compact and Hausdorff, the closed subset $C$ is compact, and the complement of the compact (hence closed) set $K$, $U$, is open. Thus:
    $$ \iota^{-1}(S(K, V)) = \{f \in H(X) \mid f(C) \subseteq U\} = S(C, U) $$
    Since $S(C, U)$ is a subbasic open set, the preimage of any subbasic open set is open, so $\iota$ is continuous.

    For $\mu(f, g) = f \circ g$,
    let $f, g \in H(X)$ such that $\mu(f, g) \in S(K, V)$. Then $(f \circ g)(K) \subseteq V$, which implies $g(K) \subseteq f^{-1}(V)$.

    The set $g(K)$ is compact (as the continuous image of a compact set) and $f^{-1}(V)$ is open. Since a compact Hausdorff space is regular, we can interpolate an open set $U$ such that:
    $$ g(K) \subseteq U \subseteq \overline{U} \subseteq f^{-1}(V) $$
    Because $X$ is compact and $\overline{U}$ is closed, $\overline{U}$ is compact.

    Define the open neighborhoods $W_f = S(\overline{U}, V)$ and $W_g = S(K, U)$. By construction, $f \in W_f$ and $g \in W_g$.
    For any $f' \in W_f$ and $g' \in W_g$, we evaluate the composition on $K$:
    $$ (f' \circ g')(K) = f'(g'(K)) \subseteq f'(U) \subseteq f'(\overline{U}) \subseteq V $$
    Thus, $\mu(W_f \times W_g) \subseteq S(K, V)$, proving that $\mu$ is continuous at $(f,g)$.

    Since both operations are continuous, $H(X)$ is a topological group.



    Now let $\alpha: G \times X \to X$ be a group action, and let $\Phi: G \to H(X)$ be the corresponding algebraic homomorphism defined by $\Phi(g)(x) = \alpha(g, x)$. We must show that $\alpha$ is continuous if and only if $\Phi$ is continuous.

    Assume $\alpha$ is continuous. To show $\Phi$ is continuous, let $S(K, V)$ be a subbasic open set in $H(X)$ and let $g \in \Phi^{-1}(S(K, V))$. By definition, $\alpha(\{g\} \times K) \subseteq V$.
    Since $\alpha$ is continuous and $K$ is compact, a standard compactness argument (the tube lemma) guarantees the existence of an open neighborhood $U$ of $g$ in $G$ such that $\alpha(U \times K) \subseteq V$.
    Therefore, for all $h \in U$, $\Phi(h)(K) \subseteq V$, which means $h \in \Phi^{-1}(S(K, V))$. Since $U \subseteq \Phi^{-1}(S(K, V))$, the preimage is open, making $\Phi$ continuous.

    Conversely assume $\Phi$ is continuous. Let $(g, x) \in G \times X$ and let $V$ be an open neighborhood of $\alpha(g, x) = \Phi(g)(x)$. Since $\Phi(g)$ is continuous, $\Phi(g)^{-1}(V)$ is an open neighborhood of $x$.
    Because $X$ is a compact Hausdorff space, it is regular. Thus, there exists an open neighborhood $W$ of $x$ such that its closure $\overline{W}$ is compact and $\overline{W} \subseteq \Phi(g)^{-1}(V)$.
    This containment implies $\Phi(g)(\overline{W}) \subseteq V$, meaning $\Phi(g) \in S(\overline{W}, V)$.
    Since $\Phi$ is continuous, $U = \Phi^{-1}(S(\overline{W}, V))$ is an open neighborhood of $g$ in $G$.
    Now, consider the open neighborhood $U \times W$ of $(g, x)$ in $G \times X$. For any $(h, y) \in U \times W$, we know $\Phi(h) \in S(\overline{W}, V)$, meaning $\Phi(h)(\overline{W}) \subseteq V$. Since $y \in W \subseteq \overline{W}$, it follows that $\alpha(h, y) = \Phi(h)(y) \in V$.
    Thus, $\alpha(U \times W) \subseteq V$, proving that the action $\alpha$ is continuous at $(g, x)$.

\end{proof}



We will also refer to a continuous action of $G$ on $X$ as a $G$-{\it flow} on $X$. If the action is understood, we will often simply use $X$ to refer to the flow.

Given a $G$-flow on $X$ and a point $x\in X$, the {\it orbit} of $x$ is the
set
\[
    G\cdot x=\{g\cdot x:g\in G\}
\]
and the {\it orbit closure} of $x$, the set
\[
    \overline{G\cdot x}.
\]
A  subset $A \subseteq X$ is \textit{$G$-invariant} if for every $g \in G$, $g \cdot A \subseteq A$.


\begin{theorem}
    The  orbit closure of $x$ is a $G$-invariant, compact subset of $X$.
\end{theorem}

\begin{proof}
    Let $A = \overline{G \cdot x}$, by definition the map $\tau_g: X \to X$ defined by $\tau_g(y) = g \cdot y$ is a homeomorphism for each $g \in G$.

    First, we note that the orbit $G \cdot x$ is itself $G$-invariant. For any $h \in G$:
    $$h \cdot (G \cdot x) = \{h \cdot (g \cdot x) : g \in G\} = \{(hg) \cdot x : g \in G\}$$
    Since $G$ is a group, the map $g \mapsto hg$ is a bijection of $G$ onto itself. Thus:
    $$h \cdot (G \cdot x) = G \cdot x$$

    For any homeomorphism $f$ and any subset $S \subseteq X$, we have $f(\overline{S}) = \overline{f(S)}$. Applying this to the map $\tau_h$:
    $$h \cdot \overline{G \cdot x} = \overline{h \cdot (G \cdot x)}$$
    Substituting the invariance of the orbit shown above:
    $$h \cdot \overline{G \cdot x} = \overline{G \cdot x}$$
    Since this equality holds for all $h \in G$, the set $\overline{G \cdot x}$ is $G$-invariant.

    The compactness follows directly from the topological properties of \(X\). By definition, the closure \(\overline{G \cdot x}\) of the orbit \(G \cdot x\) is a closed subset of \(X\). In the definition of a \(G\)-flow, the phase space \(X\) is assumed to be compact. Every closed subset of a compact space is itself compact. Consequently, \(\overline{G \cdot x}\) is compact.
\end{proof}



In general, a (non-$\emptyset$) compact, $G$-invariant subset $Y\subseteq X$ defines a {\it subflow} by restricting the $G$-action to $Y$.
\begin{mynote}
    \begin{proof}
        does it need a proof?

    \end{proof}
\end{mynote}

A $G$-flow on $X$ is {\it minimal} if
it contains no proper subflows, i.e., there is no (non-$\emptyset$) compact $G$-invariant set other than $X$.
Thus

\begin{theorem}
    $X$ is minimal iff every orbit is dense.
\end{theorem}

\begin{proof}
    Assume first that $X$ is minimal, and let $x \in X$. The orbit closure
    \[
        \overline{Gx}
    \]
    is a nonempty closed $G$-invariant subset of $X$, hence a subflow. By minimality, it must equal $X$. Therefore $Gx$ is dense in $X$.

    Conversely, suppose that every orbit is dense in $X$, and let $Y \subseteq X$ be a nonempty compact $G$-invariant subset. Pick any $y \in Y$. Since $Y$ is $G$-invariant, we have $Gy \subseteq Y$, and since $Y$ is closed, it follows that
    \[
        \overline{Gy} \subseteq Y.
    \]
    But by assumption the orbit $Gy$ is dense in $X$, so $\overline{Gy}=X$. Hence $X \subseteq Y$, and therefore $Y=X$.

    Thus $X$ has no proper nonempty compact $G$-invariant subsets, so $X$ is minimal.
\end{proof}

A simple application of Zorn's Lemma shows that every $G$-flow $X$ contains a minimal subflow $Y\subseteq X$.

\begin{proof}
    Let
    \[
        \mathcal{F}=\{Y\subseteq X : Y \neq \emptyset,\ Y \text{ is compact and } G\text{-invariant}\},
    \]
    partially ordered by inclusion. Since $X \in \mathcal{F}$, the set $\mathcal{F}$ is nonempty.

    We claim that every chain in $\mathcal{F}$ has a lower bound in $\mathcal{F}$. Let $\mathcal{C}\subseteq \mathcal{F}$ be a chain. Since $\mathcal{C}$ is linearly ordered by inclusion, it has the finite intersection property: for any $Y_1,\dots,Y_n\in\mathcal{C}$, one of them is contained in all the others, hence
    \[
        Y_1\cap\cdots\cap Y_n \neq \emptyset.
    \]
    Each member of $\mathcal{C}$ is closed in the compact space $X$, so compactness implies that
    \[
        Y_*:=\bigcap_{Y\in\mathcal{C}} Y \neq \emptyset.
    \]
    Moreover, $Y_*$ is compact, being a closed subset of the compact space $X$, and it is $G$-invariant, since an intersection of $G$-invariant sets is again $G$-invariant. Thus $Y_* \in \mathcal{F}$, and clearly $Y_* \subseteq Y$ for every $Y\in\mathcal{C}$. So $Y_*$ is a lower bound for $\mathcal{C}$ in $\mathcal{F}$.

    By Zorn's Lemma, $\mathcal{F}$ has a minimal element $Y$. By definition, $Y$ is a nonempty compact $G$-invariant subset of $X$ with no proper nonempty compact $G$-invariant subset. In other words, $Y$ is a minimal subflow of $X$.
\end{proof}


Among minimal flows of a given group $G$, there is a largest (universal) one, called the {\it universal minimal flow}.  To define this, we first need the concept of homomorphism of $G$-flows.  Let $X,Y$ be two $G$-flows. A {\it homomorphism} of the $G$-flow $X$ to the $G$-flow $Y$ is a continuous map $\pi :X\rightarrow Y$, which is also a $G$-{\it map}, i.e.,
\[
    \pi (g\cdot x)=g\cdot \pi (x),\ \ \ x\in X, g\in G.
\]
Notice that if $Y$ is minimal, then any homomorphism of $X$ into $Y$ is
surjective.
\begin{proof}
    Let $\pi:X\to Y$ be a homomorphism of $G$-flows, and assume that $Y$ is minimal. Since $X$ is nonempty and $\pi$ is continuous, the image $\pi(X)$ is a nonempty compact subset of $Y$. We claim that $\pi(X)$ is $G$-invariant.

    Indeed, if $y\in \pi(X)$, then $y=\pi(x)$ for some $x\in X$. For any $g\in G$, we have
    \[
        g\cdot y=g\cdot \pi(x)=\pi(g\cdot x)\in \pi(X),
    \]
    since $\pi$ is a $G$-map. Thus $\pi(X)$ is $G$-invariant.

    Therefore $\pi(X)$ is a nonempty compact $G$-invariant subset of $Y$. By minimality of $Y$, there are no proper subsets of this kind, so necessarily
    \[
        \pi(X)=Y.
    \]
    Hence $\pi$ is surjective.
\end{proof}
An {\it isomorphism} of $X$ to $Y$ is a bijective homomorphism
$\pi :X\rightarrow Y$ (notice then that $\pi^{-1}$ is also a homomorphism).
We now have the following basic fact in topological dynamics.

\begin{theorem}
    \label{1.1}
    Given a topological group $G$, there is a
    minimal $G$-flow $M(G)$ with the following property:  For any minimal
    $G$-flow $X$ there is a homomorphism $\pi :M(G)\rightarrow X$.  Moreover,
    $M(G)$ is uniquely determined up to isomorphism by this property.
\end{theorem}

\begin{proof}
    See \ref{thm:universal_minimal_flow_existence_uniqueness}.
\end{proof}


The space $M(G)$ is called the {\it universal minimal flow} of $G$.  In order
to get an intuition about this space, we will discuss a standard way of
looking at it.  For that we first need to discuss the concept of a {\it
        pointed} $G$-{\it flow} or {\it ambit}.

Let $G$ be a topological group.  A $G$-{\it ambit} is a $G$-flow $X$ with a
distinguished point $x_0\in X$, whose orbit is dense in $X$.  We often
abbreviate this by $(X,x_0)$.  A {\it homomorphism of} $G$-{\it ambits}
$(X,x_0), (Y,y_0)$ is a homomorphism $\pi :X\to Y$ of the $G$-flows
such that $\pi (x_0)=y_0$. If such a homomorphism exists, it is clearly
unique.  Similarly we define the concept of isomorphism of $G$-ambits.

It is another basic fact of topological dynamics that there is again
a largest (universal) $G$-ambit.

\begin{theorem}
    \label{1.2}
    Given a topological group $G$, there is a
    $G$-ambit ($S(G), s_0$) with the following property:  For any $G$-ambit
    $(X,x_0)$ there is a homomorphism of $(S(G), s_0)$ to $(X,x_0)$.  Moreover,
    $(S(G), s_0)$ is uniquely determined up to isomorphism by this property.
    \qed
\end{theorem}

\begin{proof}
    See \ref{thm:universal_T_ambit_existence_uniqueness}.
\end{proof}

The space $(S(G),s_0)$, often simply written as $S(G)$, is called the {\it
        greatest G-ambit}.

The uniqueness part of the preceding result is obvious from the definitions.


\begin{mynote}
    Shall I continue from here?
\end{mynote}

To establish existence, we will describe a particularly useful way of
constructing the greatest ambit.

Consider the space $RUC^b(G)$ of bounded right-uniformly continuous functions $x:G\rightarrow\bbC$.
Recall that $x:G\rightarrow\bbC$ is {\it right-uniformly continuous} if for
each $\epsilon >0$ there is a nbhd $V$ of the identity $1_G$ of $G$ so that
\[
    gh^{-1}\in V\Rightarrow |x(g)-x(h)|<\epsilon .
\]
If $G$ is metrizable, with right-invariant compatible metric $d_r$, then
right-uniformly continuous means uniformly continuous with respect to $d_r$:
\[
    \forall\epsilon\exists\delta (d_r(g,h)<\delta\Rightarrow |x(g)-x(h)|<
    \epsilon ).
\]
Under pointwise addition, multiplication and conjugation, and with the
sup norm $\Vert x\Vert_\infty =\sup\{|x(g)|:g\in G\}, RUC^b(G)$ is an
abelian $C^*$-algebra which is unital (with multiplicative identity
the constant 1 function).

\begin{mynote}
    \begin{proposition*}
        $RUC^b(G)$ is a unital abelian $C^*$-algebra.
    \end{proposition*}

    \begin{proof}
        Recall that
        \[
            RUC^b(G)=\{x:G\to\mathbb C:\ x \text{ is bounded and right-uniformly continuous}\},
        \]
        with pointwise operations, involution $x^*(g)=\overline{x(g)}$, and the norm
        \[
            \|x\|_\infty=\sup_{g\in G}|x(g)|.
        \]

        \medskip
        \noindent\textbf{1. $RUC^b(G)$ is a complex vector space.}
        Let $x,y\in RUC^b(G)$ and $\lambda\in\mathbb C$.
        Boundedness and right-uniform continuity are preserved under pointwise
        addition and scalar multiplication, using the usual $\varepsilon/2$ arguments.
        Thus $x+y,\ \lambda x\in RUC^b(G)$.

        The vector space axioms hold pointwise:
        \[
            (x+y)(g)=x(g)+y(g),\qquad (\lambda x)(g)=\lambda x(g),
        \]
        so $RUC^b(G)$ is a complex vector space.

        \medskip
        \noindent\textbf{2. $RUC^b(G)$ is an abelian algebra.}
        Define multiplication by
        \[
            (xy)(g)=x(g)y(g).
        \]
        This multiplication is associative and commutative pointwise, and bilinear:
        \[
            (x+y)z = xz + yz,\qquad x(y+z)=xy+xz,\qquad (\lambda x)y=\lambda(xy).
        \]

        We also check that $xy$ is right-uniformly continuous.
        Let $x,y\in RUC^b(G)$ and fix $\varepsilon>0$.
        Since $x$ and $y$ are bounded, let $M=\max(\|x\|_\infty,\|y\|_\infty)$.
        Choose neighborhoods $V_x,V_y$ of the identity such that
        \[
            gh^{-1}\in V_x\Rightarrow |x(g)-x(h)|<\frac{\varepsilon}{2M},\qquad
            gh^{-1}\in V_y\Rightarrow |y(g)-y(h)|<\frac{\varepsilon}{2M}.
        \]
        Let $V=V_x\cap V_y$.  For $gh^{-1}\in V$,
        \begin{align*}
            |x(g)y(g)-x(h)y(h)|
             & \le |x(g)|\,|y(g)-y(h)| + |y(h)|\,|x(g)-x(h)|                 \\
             & < M\cdot\frac{\varepsilon}{2M} + M\cdot\frac{\varepsilon}{2M}
            = \varepsilon.
        \end{align*}
        Thus $xy\in RUC^b(G)$.
        Hence $RUC^b(G)$ is an abelian algebra.

        \medskip
        \noindent\textbf{3. Submultiplicativity of the sup-norm.}
        For $x,y\in RUC^b(G)$,
        \[
            \|xy\|_\infty
            =\sup_{g\in G}|x(g)y(g)|
            \le \sup_{g}|x(g)|\cdot\sup_{g}|y(g)|
            =\|x\|_\infty\|y\|_\infty.
        \]
        Hence $(RUC^b(G),\|\cdot\|_\infty)$ is a normed algebra.

        \medskip
        \noindent\textbf{4. Completeness (Banach algebra).}
        Let $(x_n)$ be a Cauchy sequence in $\|\cdot\|_\infty$.
        Since the space of bounded functions on $G$ with the sup-norm is complete,
        $x_n$ converges uniformly to a bounded function $x:G\to\mathbb C$.

        We show $x$ is right-uniformly continuous.
        Let $\varepsilon>0$.  Choose $N$ such that $\|x_n-x\|_\infty<\varepsilon/3$.
        Since $x_n\in RUC^b(G)$, choose a neighborhood $V$ of $1_G$ such that
        \[
            gh^{-1}\in V \Rightarrow |x_n(g)-x_n(h)|<\varepsilon/3.
        \]
        Then for such $g,h$,
        \[
            |x(g)-x(h)|
            \le |x(g)-x_n(g)| + |x_n(g)-x_n(h)| + |x_n(h)-x(h)|
            < \varepsilon.
        \]
        Thus $x\in RUC^b(G)$, so $RUC^b(G)$ is complete.

        \medskip
        \noindent\textbf{5. Involution.}
        Define $x^*(g)=\overline{x(g)}$.  This satisfies:
        \[
            (x^*)^*=x,\qquad (x+y)^*=x^*+y^*,\qquad (\lambda x)^*=\overline{\lambda}x^*,\qquad
            (xy)^*=y^*x^*,
        \]
        all pointwise.  Moreover,
        \[
            |x^*(g)-x^*(h)| = |\overline{x(g)}-\overline{x(h)}| = |x(g)-x(h)|,
        \]
        so right-uniform continuity is preserved.
        Thus $RUC^b(G)$ is a $*$-algebra.

        \medskip
        \noindent\textbf{6. The $C^*$-identity.}
        For every $x\in RUC^b(G)$,
        \[
            \|x^*x\|_\infty
            =\sup_{g\in G}|\overline{x(g)}\,x(g)|
            =\sup_{g}|x(g)|^2
            =\bigl(\sup_{g}|x(g)|\bigr)^2
            =\|x\|_\infty^2.
        \]
        Thus the $C^*$-identity holds.

        \medskip
        \noindent\textbf{7. Unit element.}
        The constant function $1(g)=1$ belongs to $RUC^b(G)$ (it is bounded and
        right-uniformly continuous).  For every $x\in RUC^b(G)$,
        \[
            (1\cdot x)(g)=1\cdot x(g)=x(g).
        \]
        Thus $RUC^b(G)$ is unital.

        \medskip
        Combining all the above, $RUC^b(G)$ is a unital abelian $C^*$-algebra.
    \end{proof}

\end{mynote}

Denote by $S(G)$ the maximal ideal space of
$RUC^b(G)$, i.e., the space of all continuous homomorphisms $\varphi:RUC^b(G)\rightarrow\bbC$.
Equipped with the topology generated by
the maps $\hat x:S(G)\rightarrow\bbC,\hat x(\varphi )=\varphi (x)$, for
$x\in RUC^b(G)$ (i.e., the smallest topology in which all $\hat x,
    x\in RUC^b(G)$, are continuous), this is a compact, Hausdorff space.
Moreover, by the Gelfand-Naimark theorem, $RUC^b(G)$ can be canonically
identified with $C(S(G))$, the $C^*$-algebra of all continuous
complex-valued functions on $S(G)$, identifying $x\in RUC^b(G)$
with $\hat x$ (see, e.g., \cite{Ru}, 11.18).

Now $G$ acts continously by $C^*$-automorphisms on $RUC^b(G)$ by
left-shift
\[
    g\cdot x(h) =x(g^{-1}h)
\]
and thus acts canonically on $S(G)$ via
\[
    g\cdot\varphi (x)=\varphi (g^{-1}\cdot x).
\]
It is also easy to check that this action is continuous, so $S(G)$ is a
$G$-flow.  We will now identify a canonical element of $S(G)$, that will
turn it into an ambit.

For each $g\in G$, let $\varphi_g\in S(G)$ be defined by $\varphi_g
    (x)=x(g),x\in RUC^b(G)$.  Then one can see that $g\mapsto\varphi_g$
is a homeomorphism of $G$ with a dense subset of $S(G)$.  For example,
when $G$ is metrizable with bounded compatible right-invariant metric
$d_r$, and $g_0\neq h_0$, then for $x(g)=d_r(g,h_0),\varphi_{g_0}(x)\neq
    \varphi_{h_0}(x)$ so $\varphi_{g_0}\neq\varphi_{h_0}$, i.e., this map is
1-1.  The verification that it is homeomorphism is straightforward.
Finally $\{\varphi_g:g\in G\}$ is dense in $S(G)$, since, otherwise,
there is $f\in C(S(G))$, so that $f=\hat x$ for some $x\in RUC^b(G)$,
with $f\neq 0$ but $f(\varphi_g)=\hat x(\varphi_g)=x(g)=0,\forall g\in G$,
which implies that $x=0$, thus $f=0$, a contradiction.

So from now on we will identify $g$ with $\varphi_g$ and think of $G$ as
a dense subset of $S(G)$.  Moreover $G$ is an invariant subset of
$S(G)$ and the restriction of the action to $G$ is simply left-translation:
$(g,h)\mapsto gh$.  We now have the following standard fact.

\begin{theorem}
    \label{1.3}
    The $G$-ambit $(S(G), 1_G)$ is the
    greatest $G$-ambit.
\end{theorem}

\begin{proof}
    Since the orbit of $1_G$ in $S(G)$ is $G$, which is dense,
    clearly $(S(G),1_G)$ is a $G$-ambit.  Consider now an arbitrary $G$-ambit
    $(X,x_0)$.  Suppose $f\in C(X)$.  Define then $f^*:G\rightarrow\bbC$ by
    \[
        f^*(g)=f(g\cdot x_0).
    \]
    We verify that $f^*\in RUC^b(G)$.  Since the action of $G$ on $X$
    is continuous, an easy compactness argument shows that given $\epsilon >0$,
    there is a nbhd $V$ of the identity of $G$ such that $g\in V$ implies
    $|f(g\cdot x)-f(x)|<\epsilon ,\forall x\in X$.  So if $gh^{-1}\in V$, then
    $|f^*(g)-f^*(h)|=|f(g\cdot x_0)-f(h\cdot x_0)|=|f(gh^{-1}\cdot (h\cdot x_0))
        -f(h\cdot x_0)|<\epsilon$, and thus $f^*\in RUC^b(G)\ (f^*$ is clearly
    bounded).

    Identifying, as usual, $RUC^b(G)$ with $C(S(G))$, the map $f\mapsto
        f^*$ is a unital $C^*$-algebra monomorphism of $C(X)$ into $C(S(G))$.
    Now it is a well known fact that a unital $C^*$-algebra monomorphism
    $\pi :C(K)\rightarrow C(L)$, where $K,L$ are (non-$\emptyset$) compact
    spaces, is of the form $\pi (f)=f\circ\Pi$ for a {\it uniquely}
    determined continuous surjection $\Pi :L\rightarrow K$ (see, e.g.,
    \cite{Co}, 2.4.3.6).  From this it also follows that if $K,L$
    are actually $G$-flows, and we let $G$ act on $C(K),C(L)$ by shift,
    $g\cdot f(x)=f(g^{-1}\cdot x)$, then if $\pi$ is a $G$-map, $\Pi$ is
    also a $G$-map.  Applying this to $f\mapsto f^*$, we see that there is a
    homomorphism of $G$-flows $\Phi :S(G)\rightarrow X$ with $f^*=f\circ\Pi$.
    It only remains to check that $\Pi (1_G)=x_0$.  But for any $f\in C(X),
        f^*(1_G)=f(x_0)=f(\Pi (1_G))$, so we must have $\Pi (1_G)=x_0$, and
    the proof is complete.
\end{proof}

Using this it is now immediate to obtain the following description of the
universal minimal flow of $G$.

\begin{corollary}
    \label{1.4}
    Let $M(G)$ be a minimal subflow of
    $S(G)$ (i.e., $M(G)$ is a minimal $G$-invariant compact subset of $S(G)$).
    Then $M(G)$ is the universal minimal flow (up to isomorphism).
\end{corollary}

\begin{proof}
    Let $X$ by any minimal $G$-flow.  Fix $x_0\in X$.  Then
    $(X,x_0)$ is a $G$-ambit, so let $\pi :(S(G), 1_G)\rightarrow (X,x_0)$
    be a homomorphism.  Then clearly the restriction of $\pi$ to $M(G)$ is also
    a homomorphism, and we are done.
\end{proof}

In particular, it follows from the uniqueness part of Theorem \ref{1.1}, that all
minimal subflows of $S(G)$ are isomorphic.  (This uniqueness part, which
is not proved here, is based on techniques from semigroup theory.)
As we will soon see, the space $M(G)$ can be extremely complicated, e.g.,
non-metrizable, even when the group $G$ is very ``small'', e.g.,
an infinite
countable discrete $G$.  However, we will verify that when $G$ is
separable, metrizable, $M(G)$ is at least an inverse limit of metrizable
$G$-flows.

Fix a topological group $G$.  An {\it inverse system} of $G$-flows consists
of a directed set $\langle I,\preceq\rangle$, a family $\{X_i\}_{i\in I}$
of $G$-flows and a family of homomorphisms $\pi_{ij}:X_j\rightarrow X_i$,
for each $i\preceq j$, such that $\pi_{ii}=$ the identity of $X_i$, and
$i\preceq j\preceq k\Rightarrow\pi_{ik}=\pi_{ij}\circ\pi_{jk}$.  The
    {\it inverse limit} $\lim_{\leftarrow}X_i$ is the $G$-flow defined as
follows:  Consider the product topological space $\prod_{i\in I}X_i$ and
let
\[
    \lim_{\leftarrow}X_i=\{\{x_i\}_i\in\prod_iX_i:\forall i\preceq j(\pi_{ij}
    (x_j)=x_i)\}.
\]
By a simple application of compactness, $\lim_{\leftarrow}X_i\neq\emptyset$ and is
clearly a compact subset of $\prod_{i\in I}X_i$.  The group $G$ acts on
$\lim_{\leftarrow}X_i$ coordinatewise:  $g\cdot\{x_i\}=\{g\cdot x_i\}$ and
this is clearly a continuous action.  Define
\[
    \pi_i:\lim_{\leftarrow}X_i\rightarrow X_i
\]
by $\pi_i(\{x_i\}_{i\in I})=x_i$.  Then $\pi_i$ is a homomorphism and
if $i\preceq j$, then $\pi_i=\pi_{ij}\circ\pi_j$.  Finally, if $X$ is a
$G$-flow and there are homomorphisms $\varphi_i:X\rightarrow X_i$ with
$i\preceq j\Rightarrow\varphi_i=\varphi_{ij}\circ\varphi_j$, then there is
a unique homomorphism $\varphi :X\rightarrow\lim_{\leftarrow}X_i$ such
that $\varphi_i=\pi_i\circ\varphi$.

Similarly we define inverse systems of $G$-ambits.  We now have the following
folklore fact.

\begin{theorem}
    \label{1.5}
    Let $G$ be a separable, metrizable group.
    Then the greatest ambit $(S(G),1_G)$ is the inverse limit of a system of
    metrizable $G$-ambits.  Similarly, the universal minimal flow is the inverse
    limit of a system of metrizable minimal $G$-flows.
\end{theorem}

\begin{proof} We will first derive the second assertion from the first.
    Suppose $\{(X_i,x^0_i)\}$ is an inverse system of metrizable $G$-ambits,
    so that $(X,x_0)=\lim_{\leftarrow}(X_i,x^0_i)$ is the greatest $G$-ambit.
    Let $\pi_i:X\rightarrow X_i$ be the corresponding homomorphism.  In
    particular, $X=\lim_{\leftarrow}X_i$ as a $G$-flow.  Now fix a minimal
    subflow $M\subseteq X$, so that, as we have seen earlier, $M$ is the
    universal minimal flow.  Put $\pi_i(M)=M_i\subseteq X_i$.  Then $M_i$ is a
    subflow of $X_i$ and it is clearly minimal and metrizable.  So it is enough
    to check that $M=\lim_{\leftarrow}M_i$, which is an easy compactness
    argument.

    For the first assertion, fix a countable dense set $D\subseteq G$.  Let
    $(I,\preceq)$ be the following directed set:  $I$ consists of all separable,
    closed, unital $G$-invariant $C^*$-subalgebras of $RUC^b(G)$.  Since a
    closed, unital $C^*$-subalgebra of $RUC^b(G)$ is $G$-invariant iff it
    is $D$-invariant, clearly $\bigcup I=RUC^b(G)$.  Let for $A,B\in I$
    \[
        A\preceq B\Leftrightarrow A\subseteq B.
    \]
    For $A\in I$ denote by $X_A$ the maximal ideal space of $A$, which is
    compact metrizable, since $A$ is separable.  Since $G$ acts continuously
    by $C^*$-automorphisms on $A$, it acts continuously on $X_A$.  Moreover,
    as before,  we can identify each $g\in G$ with an element of $X_A$, so that
    $G$ can be thought as a dense invariant subset of $X_A$ and the $G$ action
    on it is by left-translation.  Thus again $(X_A,1_G)$ is a metrizable
    $G$-ambit.  We view as usual $A$ as identified with $C(X_A)$ via $x\mapsto
        \hat x_A$.  When $A\preceq B$, the identity is an injective unital
    $C^*$-homomorphism from $A$ to $B$, so there is a unique surjection
    $\pi_{AB}:X_B\rightarrow X_A$ such that for $x\in A$, $\hat x_B=\hat x_A
        \circ\pi_{AB}$, therefore for $\varphi\in X_B, \hat x_A(\pi_{AB}(\varphi
        ))=\hat x_B(\varphi )$ or $\pi_{AB}(\varphi )(x)=\varphi (x)$, i.e.,
    $\pi_{AB}(\varphi )=\varphi |A$.  Similarly, there is a surjection
    $\pi_A:S(G)\rightarrow X_A$ given by $\pi_A(\varphi )=\varphi |A$ for
    each $A\in I$, so that $\pi_A=\pi_{AB}\circ\pi_B$ for $A\preceq B$.
    Moreover, $\pi_A (1_G)= \pi_{AB} (1_G) = 1_G$.  Thus
    $\varphi\mapsto (\pi_A(\varphi))_{A\in I}$
    is a homomorphism from $(S(G),1_G)$ to $\lim_{\leftarrow}(X_A,1_G)$.
    Also given $(\varphi_A)\in\lim_{\leftarrow}X_A$, we have, for
    $A\preceq B$, that $\varphi_A=\varphi_B|A$, so that there is a unique
    $\varphi\in S(G)$ with $\pi_A(\varphi )=\varphi_A$.  Thus $\varphi\mapsto
        (\pi_A(\varphi ))_{A\in I}$ is an isomorphism of the $G$-ambit
    $(S(G),1_G)$ and $\lim_{\leftarrow} (X_A,1_G)$.
\end{proof}


\section{B}
We will now discuss the case of infinite countable discrete
groups $G$ and see that in this case $M(G)$ is an extremely large space,
in particular it is not metrizable.

For an infinite countable discrete $G$, it is clear that $RUC^b(G)$ is
identical with $\ell^\infty (G)$, the $C^*$-algebra of bounded complex
functions on $G$ with the supremum norm.  It is then easy to see that
$S(G)$ is identical with $\beta G$, the space of ultrafilters on $G$ with
the topology whose basis consists of the sets $\hat A=\{U\in\beta G:
    A\in U\}$, for $\emptyset\neq A\subseteq G$.  This is a non-metrizable,
compact, Hausdorff space.  The action of $G$ on $S(G)$ is given by
\[
    A\in g\cdot U\Leftrightarrow g^{-1}A\in U
\]
for $A\subseteq G$.

The copy of $G$ in $S(G)$ consists simply of the principal ultrafilters.
So the distinguished point is the principal ultrafilter on $1_G$.  Finally, $M(G)$ is any minimal subflow
of $\beta G$.

First we point out that $G$ acts {\it freely} on $S(G)$, i.e., $g\cdot
    x\neq x,\ \forall g\neq 1_G,\ x\in S(G)$ (this is due to Ellis \cite{E60}).
For that it is of course enough to find a free $G$-flow $X$
(since $S(G)$ can be homomorphically mapped to $X$).  Again it is enough
to find for each $g\in G,\ g\neq 1_G$, a $G$-flow $X_g$, such
that $g\cdot x\neq x,\ \forall x\in X_g$.  Because then $X=\prod_{g\in G
        \setminus\{1_G\}}X_g$ with the action $h\cdot (x_g)=(h\cdot x_g)$ works.
Consider the shift action of $G$ on $3^G,\ g\cdot p(h)=p(g^{-1}h)$.
Suppose we can find $p_g\in 3^G$ such that $(*):\forall h\in G(g\cdot
    (h\cdot p_g)(1_G)\neq (h\cdot p_g)(1_G))$.  Then we can take $X_g$ to be
the orbit closure of $p_g$.  Now $(*)$ is equivalent to:  $\forall h\in G
    (p_g(h^{-1}g^{-1})\neq p_g(h^{-1}))$ or $\forall h(p_g(h)\neq p_g(hg))$.
Consider a coset $h\langle g\rangle$ of the cyclic subgroup $\langle g
    \rangle$.  Define $p_g$ on $h\langle g\rangle$, so that $p_g(hg^k)\neq
    p_g(hg^{k+1})$, for any $k\in\bbZ$.  (We used $3^G$ instead of $2^G$
to take care of the case when $g$ has finite order.)  This clearly works.

The space $M(G)$ is quite big, see e.g., the references in \cite{dV},
p. 391, {\bf 11}.  Let us verify for instance that it is not metrizable.
The space $M(G)$ is a closed subset of $\beta G$.  If it was metrizable
and infinite, it would have a non-trivial convergent sequence, which is
impossible in the extremally disconnected space $\beta G$ (see \cite{Eng},
Ex. 6.2.G(a) on p. 456).  So if $M(G)$ is metrizable, it has to be
finite, contradicting the fact that $G$ acts freely on $M(G)$.
\medskip

\section{C}
More generally, Veech \cite{V} has shown that when $G$
is locally compact, then
$G$ acts freely on $S(G)$ and thus on $M(G)$.  For a simpler version of
the proof see Pym \cite{Py}.  See also \cite{AS} for the second countable
case.  In an Appendix to this paper, we also give a new proof of Veech's
Theorem.  Note that Veech's Theorem implies that if $G$
is second countable, then $G$
admits a free metrizable $G$-flow.  To see this, notice that it is enough
to find for each $1_G\neq g\in G$ a metrizable $G$-flow $X_g$ with
$g\cdot x\neq x$, if $x\in X_g$.  Indeed, if we can do that, then, by
compactness, there is an open nbhd $V_g$ of $g$ with $1_G\not\in V_g$ and $h\in
    V_g\Rightarrow (h\cdot x\neq x,\ \forall x\in X_g)$.  Find now $g_0,g_1,
    \dots \in G\setminus\{1_G\}$ so that $\{V_{g_n}\}_{n\in\bbN}$ is an open
cover of $G\setminus\{1_G\}$.  Then $X=\prod_nX_{g_n}$ with the
coordinatewise action is a free metrizable $G$-flow. (This argument comes
from \cite{AS}.)
So fix $g\in G,\ g\in 1_G$, in order to find $X_g$.  Write $S(G)=
    \lim_{\leftarrow}X_i,\ X_i$ a metrizable $G$-flow.  If none of the
$X_i$ can be $X_g,\ \{x_i\in X_i:g\cdot x_i=x_i\}=Y_i$ is non-$\emptyset$,
and if $\pi_{ij}:X_j\rightarrow X_i$, for $i\preceq j$, are the
corresponding homomorphisms, then $\pi_{ij}(Y_j)\subseteq Y_i$, so the
inverse limit of $(Y_i,\pi_{ij}|Y_j)$ is non-$\emptyset$ and thus there
is $\{y_i\}\in S(G)$ with $g\cdot \{y_i\}=\{y_i\}$, a contradiction, as $G$
acts freely on $S(G)$.

We also show in an Appendix that when $G$ is non-compact, locally compact,
$M(G)$ is non-metrizable.  Stronger results in
special cases were obtained in Turek \cite{Tu}, Lau-Milner-Pym \cite{LMP}.  Of
course when $G$ is compact, $M(G)$ is $G$ itself with the left-translation
action.
\medskip

\section{D}
Rather remarkably, there are nontrivial groups $G$ for which $M(G)$
trivializes, i.e., consists of a single point.  Such groups are called
extremely amenable.  Thus a topological group is {\it extremely
        amenable} if any $G$-flow $X$ has a fixed point, i.e., there is $x\in X$
with $g\cdot x = x,\ \forall g\in G$.  (For this reason, sometimes
extremely amenable groups are described as groups having the {\it fixed
        point on compacta property}.)  By Veech's Theorem nontrivial such
groups cannot be
locally compact.  As it turned out, a number of important, non-locally
compact Polish groups are extremely amenable.  Among them are: the unitary
group $U(H)$ of the infinite dimensional separable Hilbert space $H$
(Gromov-Milman \cite{GroM}); $L(I,\mathbb T )$, the group of measurable maps from $I=
    [0,1]$ to $\mathbb T$, with pointwise multiplication, and the topology of
convergence in measure (Furstenberg-Weiss, Glasner \cite{Gl98}); $H_+(I)$ and
$H_+(\bbR )$, the groups of orientation preserving homeomorphisms of
$I$ and $\bbR$, with the compact-open topology (Pestov \cite{P98a}); Aut$(I,
    \lambda )$ (resp., Aut$^*(I,\lambda ))$, the groups of measure preserving
(resp., measure-class preserving) automorphisms of Lebesgue measure
$\lambda$ on $I$, with the weak topology (Giordano-Pestov \cite{GiP});
Iso($\bf U$), the group of isometries of the Urysohn space, with the
pointwise convergence topology (Pestov \cite{P02})), and Aut$(\langle\bbQ ,
    <\rangle )$, the group of automorphisms of the rationals with the usual
ordering, with the pointwise convergence topology (Pestov \cite{P98a}).  For
more about extreme amenability, see also Pestov \cite{P99},
\cite{P02a}, \cite{P02b} and Uspenskij \cite{Usp02}.

In case the group $G$ is separable metrizable, we can restate the definition
of extreme amenability in terms of metrizable flows only.  In
other words, a separable metrizable group $G$ is extremely amenable iff
every metrizable $G$-flow has a fixed point.  Indeed, if every
metrizable $G$-flow has a fixed point, every minimal metrizable $G$-flow
is a singleton, and thus so is $M(G)$, being an inverse limit of such $G$-flows.
\medskip


\section{E}
Except for the case of compact metrizable or extremely amenable groups, there were very
few cases where the universal minimal flow $M(G)$ was known to be metrizable.
Pestov \cite{P98a} first computed that for the group $H_+(\mathbb T )$ of
orientation-preserving homeomorphisms of the circle $\mathbb T$, the canonical
evaluation action on $\mathbb T$ is the universal minimal flow.  Then
Glasner-Weiss \cite{GlW02} computed the universal minimal flow of $S_\infty$,
the group of permutations of $\bbN$ with the pointwise convergence topology.
It turns out to be the canonical action on the compact, metrizable space of
linear orderings on $\bbN$.  Finally Glasner-Weiss \cite{GlW03} computed the universal minimal flow of the
group $H(2^\bbN )$ of the homeomorphisms of the Cantor space $2^\bbN$.
It is the action of $H(2^\bbN )$ on the space of maximal chains of compact
subsets of $2^\bbN$, invented by Uspenskij \cite{Usp00}.

Let us point out here that only one of the following is possible:

(i) $M(G)$ is finite.

(ii) $M(G)$ is perfect (i.e., has no isolated points), and thus card$(M(G))=
    2^{\aleph_0}$, if $M(G)$ is metrizable.

To see this, note that if $M(G)$ has an isolated point $x_0$, then, as
every orbit of $M(G)$ is dense, $x_0$ belongs to every orbit, thus there is
only one orbit in $M(G)$, and so every point of $M(G)$ is isolated.
Since $M(G)$ is compact, $M(G)$ is finite.

We have already discussed examples of metrizable $M(G)$ which
consist of exactly one point or are perfect.  It is easy to see that also
any finite cardinality for $M(G)$ is possible.  Because if $H$ is extremely
amenable, and $G=H\times\bbZ_n$, then the obvious action of $G$ on $\bbZ_n$
is the universal minimal flow of $G$, thus card$(M(G))=n$.

Our goal in this paper is to study extreme amenability and universal
minimal flows of closed subgroups of $S_\infty$, i.e., automorphism
groups of countable structures.  In particular, we find new examples of
extremely amenable groups and also new cases where the universal minimal
flow is metrizable and can be computed.
